\subsection{History and relation to earlier work}
\label{sec:history}

\Cref{thm:main} answers positively the degreewise finite-generation
question stated by Hovey--Strickland \citep[Problem~16.2]{hoveyStrickland1999}
and in Hovey's Morava $K$- and $E$-theory problem list
\citep[Problem~1]{hoveyProblems}.
In Li's terminology, the question asks whether the local sphere is of
\emph{finite type}: its homotopy groups are finitely generated over $\Zp$
in every integer degree \citep[Introduction]{liFiniteType2021}.

The classical height-one calculation expresses the local sphere as the
fiber of an Adams operation on completed $K$-theory
\citep[\S\S1--2]{hopkinsKOne}; we recall the all-degree answer at odd
primes in \cref{prop:sphere-height-one}.
At height two and $p\geq5$, Hovey--Strickland's analysis of Shimomura's
calculations proves that $\pi_aL_{K(2)}F$ is finite for every finite
torsion spectrum $F$ and every integer $a$
\citep[\S15.2 and Theorem~15.1]{hoveyStrickland1999}.
Applied to the Moore spectrum, this supplies the mod-$p$ finiteness
criterion for the sphere in that range. The proof here establishes
that criterion uniformly in the prime and height.

Devinatz approached integral finiteness through continuous homotopy
fixed points for smaller subgroups of the Morava stabilizer. He used
\emph{essentially finite rank}, a condition of finite generation over
the periodicity ring after quotienting by the closure of periodicity
torsion \citep[\S4]{devinatzHomotopy2007}. For Morava $E$-theory $E_n$
and certain closed procyclic stabilizer subgroups $H$, he proved
this property for $\pi_*(E_n^{hH}\wedge F)$ under chromatic-acyclicity
and annihilation hypotheses on the finite spectrum $F$
\citep{devinatzFiniteness2008}.
Such intermediate fixed-point spectra require their own finiteness
condition; these results are distinct from degreewise finite
generation for the full local sphere.

Barthel--Schlank--Stapleton--Weinstein determined the rational answer
at every prime and height \citep[Theorem~A]{bssw2024}:
\[
 \Q\otimes_{\Z}\pi_*L_{K(n)}S
 \cong \Lambda_{\Qp}(\zeta_1,\ldots,\zeta_n),
 \qquad |\zeta_i|=1-2i.
\]
Their coefficient comparison also gives a split injection from
Witt-vector cohomology into degree-zero Lubin--Tate coefficient
cohomology whose complement in each cohomological degree is killed
by a power of $p$
\citep[Theorem~B]{bssw2024}.
These conclusions do not bound the size of integral torsion: the compact
module $\prod_{j\geq1}\Fp$ has zero rationalization and exponent $p$,
but is not finitely generated over $\Zp$.
The present theorem proves that the torsion subgroup is finite in
each integer homotopy degree.
Combined with the rational calculation, it determines the free ranks
in \cref{cor:sphere-ranks}.

Continuous Morava descent links coefficient cohomology to sphere
homotopy \citep[Theorem~1 and Appendix~A]{devinatzHopkins2004}.
The Hopkins--Ravenel smash product theorem, in the descendability
formulation of Mathew, gives a horizontal vanishing line at a finite
page \citep[Corollary~4.4, Theorem~4.18 and Proposition~10.10]{mathew2016}.
Heard states the resulting strongly convergent local Adams spectral
sequence at every prime and height, with a finite-page bound
independent of the input spectrum
\citep[Theorem~4.6 and Proposition~4.13]{heard2023}.
Consequently, finite coefficient cohomology for the Moore spectrum
is enough to obtain finite homotopy in each degree. Establishing that
coefficient finiteness is the main task of this paper.

Our geometric argument places Lubin--Tate deformation theory and
Strauch's torsion strata \citep{lubinTate1966,strauch2007} in the
geometry of the Fargues--Fontaine curve \citep{farguesFontaine2018}.
Scholze--Weinstein relate $p$-divisible groups and their universal
covers to sections of associated bundles \citep{scholzeWeinstein2013}; the relative bundle
and cohomology results used here are those of Fargues--Scholze
\citep{farguesScholze2021}.
The Banach--Colmez theory initiated by Colmez \citep{colmez2002}
was related to coherent sheaves on the curve by Le Bras
\citep{lebras2018}. Ansch\"utz--Le~Bras's relative full-faithfulness
theorem supplies the passage from the compatible linear morphisms
of section sheaves constructed below to bundle maps in families
\citep[Corollary~3.11]{anschutzLeBras2025}.
Within this framework, \cref{sec:completed-tower,sec:kummer,sec:independent-tuples}
construct the marked quotient on each height stratum and its
compact-support comparisons, retaining the marking actions and the
directions of the support limits.

The passage from geometric supports back to compact coefficients
also uses analytic-group duality. Lazard's theory and the formulation
of Beaudry--Goerss--Hopkins--Stojanoska provide the uniform-subgroup
resolutions and their orientation data \citep{lazard1965,bghs}.
The tensor-induction method, recalled by Symonds in its profinite
form, passes from an open subgroup to a complex with finite isotropy
\citep[Theorem~5.2]{symonds2007}.
In \cref{sec:continuous}, trace-one elements for these finite
stabilizers replace averaging by their orders. This permits the
complete-step coefficient comparisons and full-group norm argument
even when the stabilizer has $p$-torsion.

The transfer argument combines classical constructions with a
uniformity problem at finite marking level. Cartier's $p$-basis
construction identifies differential forms with Frobenius-twisted
de Rham cohomology \citep[\S3, Theorem~1]{cartier1957}; we use the
inverse-linear operator in the conventions of Blickle--B\"ockle
\citep[Definition~2.1 and Example~2.23]{blickleBockle2011}.
Lipman's formal residue transformation law gives the finite-root
comparison \citep[Lemmas~(7.2)(b) and~(7.3)]{lipmanResidues1984}.
The additional work is to control the poles of all transfer iterates
and to choose one finite marking level for all sufficiently high
transfers. These are the steps that return the geometric finiteness
statement to the compact coefficient modules
(\cref{sec:transfers,sec:laurent,sec:cutoff,sec:pro-transfer}).

Degreewise finite generation does not extend to every invertible
$K(n)$-local spectrum.
At height two and $p\geq5$, Hovey--Strickland's analysis of
$p$-adically graded spheres gives invertible spectra that are not of
finite type \citep[Theorem~15.1]{hoveyStrickland1999}; Li classifies
the finite-type invertibles in this range
\citep[Theorem~4.22]{liFiniteType2021}.
Thus arbitrary invertible or dualizable local spectra cannot replace
the finite spectra in \cref{thm:main}.
The proof makes no chromatic splitting assumption, and it does not
compute the finite torsion summands or the chromatic fracture maps.
